\section{来源审计：本讲不是“把 Transformer 搬到图像”}

这场讲座常被压缩成一句话：把图像切成 patch，再送入 Transformer。这个摘要漏掉了真正的方法论。Lucas Beyer 先花近一半时间讨论如何定义“通用视觉表示”、如何用 VTAB 测量快速迁移、为何 self-supervision 并未自动复制 NLP 的成功、以及 Big Transfer 如何把数据、模型容量、训练耐心与去重连接起来；ViT 只是在这条证据链上进一步削弱视觉归纳偏置的实验。

原视频没有公开独立 slide PDF。本文从 Stanford Online 官方录像中恢复 36 张 distinct teaching slides，并把 bumper、纯分隔页、重复的 progressive build、Q\&A 中反复跳回的旧页面与黑场标为 intentional omissions。课程发生于 2021 年，正文保留当时的研究语境；后续发展只在“延伸”中明确标注，不倒写进讲者原意。

\lecturefigure{slide-01-title.jpg}{Lucas Beyer 的讲座标题把重点放在 large-scale visual representation learning，而不是单一架构。}{官方视频 00:28。}

\begin{knowledgebox}{读图：标题中的三个限定词}
\textbf{Large-scale} 说明结论依赖数据、模型与计算规模；\textbf{visual representation} 说明目标是可迁移的中间表示，不只是某个分类器；\textbf{learning} 说明研究问题是“哪些结构可由数据学出”，而不是预设 Transformer 天然适合所有视觉任务。后文所有曲线都应带着这三个限定词阅读。
\end{knowledgebox}

\begin{importantbox}{本讲的四个闭环}
\begin{enumerate}
  \item \textbf{Measurement loop}：通用表示 $\rightarrow$ 少量样本适配 $\rightarrow$ 多任务测量 $\rightarrow$ 修正目标；
  \item \textbf{Pretraining loop}：数据规模 $\times$ 模型容量 $\times$ recipe $\rightarrow$ 可迁移表示；
  \item \textbf{Architecture loop}：patch tokenization $\rightarrow$ sequence model $\rightarrow$ learned spatial structure；
  \item \textbf{Scaling loop}：compute budget $\rightarrow$ model shape / schedule / regularization $\rightarrow$ sample efficiency。
\end{enumerate}
\end{importantbox}

\begin{warningbox}{历史结果不能脱离 protocol 复述}
“ViT 超过 ResNet”不是无条件命题。它至少依赖 pretraining dataset、训练步数、augmentation、regularization、参数/compute budget、patch size、fine-tuning protocol 与测试集。省略这些条件，会把一条实验曲线误写成架构定律。
\end{warningbox}

\subsection{本章小结}

本讲研究的不是一个 block，而是如何在规模化实验中逐步移除视觉先验，并用迁移、鲁棒性、效率和泄漏审计验证所得表示。36 张 slide 是证据骨架，教师口头解释负责连接“为什么要做”和“结果不能说明什么”。

